\documentclass[11pt,a4paper]{article}
\usepackage[T1]{fontenc}
\usepackage{lmodern,amsmath,amssymb,amsthm,booktabs,array}
\usepackage[margin=1in]{geometry}
\usepackage[colorlinks=true,allcolors=blue]{hyperref}
\usepackage{xurl}
\hypersetup{pdftitle={Extensions of Ma--Zhang's negative answer to Delgado, S\"uer and Ye\c{s}il's symmetrization question},pdfauthor={Haobo Ma; Wenlin Zhang},pdfsubject={Numerical semigroups and row-deletion symmetrization; CC BY 4.0}}

\newtheorem{theorem}{Theorem}
\newtheorem{lemma}[theorem]{Lemma}
\newtheorem{proposition}[theorem]{Proposition}
\newtheorem{corollary}[theorem]{Corollary}
\theoremstyle{definition}
\newtheorem{definition}[theorem]{Definition}
\newtheorem{conjecture}[theorem]{Conjecture}
\newtheorem{example}[theorem]{Example}
\theoremstyle{remark}
\newtheorem{remark}[theorem]{Remark}

\title{Extensions of Ma--Zhang's negative answer to Delgado, S\"uer and Ye\c{s}il's symmetrization question}
\author{Haobo Ma\\\small ChronoAI Pte Ltd, Singapore
\and Wenlin Zhang\\\small National University of Singapore;\\\small ChronoAI Pte Ltd, Singapore}
\date{25 November 2026}


\usepackage{tabularx}
\usepackage{needspace}
\setlength{\emergencystretch}{2em}
\newcommand{\N}{\mathbb{N}_0}
\newcommand{\PF}{\operatorname{PF}}
\newcommand{\AS}{\mathcal{AS}}
\newcommand{\im}{\operatorname{im}}
\newcommand{\Ap}{\operatorname{Ap}}
\begin{document}
\maketitle
\begin{abstract}
We extend Ma and Zhang's negative answer (2026, \cite[Section~2]{MZ}) to
the surjectivity question in Delgado, S\"uer and Ye\c{s}il,
\emph{A Symmetrization for Almost Symmetric Numerical Semigroups}
(arXiv:2611.14268v1, Section~5, p.~9, final paragraph):
the row-deletion map omits not only targets containing $1$, but also
some reduced symmetric targets whose diagrams have no full cell below
the diagonal.
The ordered output elements satisfy $b_i\geq2i$; equality at $b_1=2$
forces every element below the conductor to be even.
The smallest excluded reduced target has Frobenius number $7$;
under the full-cell restriction the minimum is $9$.
We give infinite excluded families, retain the exact ordinary fiber and
fixed-point characterization, and distinguish fixed-source image bounds
from fixed-target counts.
For symmetric targets of multiplicity three and Frobenius number $2h-1$,
there is one image for $h=3,4,5$ and exactly two for every $h\geq6$.
We give both families and explicit sources. At arbitrary output
multiplicity $r<h$, source multiplicity is at most $2r-2$ and source
Frobenius number is at most $2h+2r-5$, giving a finite membership test.
\end{abstract}

\section{The question and the row-deletion map}
Delgado, S\"uer and Ye\c{s}il construct a map $\Phi$ from almost symmetric
numerical semigroups to symmetric numerical sets \cite[Theorem~4.2]{DSY}.
They ask whether $\Phi$ is surjective \cite[Section~5, p.~9, final paragraph]{DSY}.
The distinction between a numerical set and a numerical semigroup is decisive:
the former may contain $1$ without being all of $\N$, whereas the latter cannot.
The omission obstruction and its consequences were established in \cite[Section~2]{MZ}. We retain those arguments and the prefix and full-cell obstructions of \cite[Sections~3--6]{MZ4}, and characterize the image at multiplicities two and three.
An earlier construction of Ye\c{s}il treats almost symmetric semigroups whose
nonmaximal pseudo-Frobenius numbers form a consecutive block
\cite[Lemma~5.7]{Y}; the general construction in \cite[Theorem~4.2]{DSY}
removes that restriction. The obstruction below does not require any
restriction on the positions of the pseudo-Frobenius numbers, or even
almost symmetry of the source.
The correspondence between omission of $1$ and a bottom row of length one
is explicit in S\"uer and Ye\c{s}il \cite[Proposition~3.5 and Corollary~3.6]{SY}.
The issue here is that this property survives the prescribed row deletion.

A \emph{numerical set} is a cofinite subset $U$ of
$\N=\{0,1,2,\ldots\}$ containing $0$. It is \emph{proper} if $U\ne\N$.
Write $G(U)=\N\setminus U$ for its nonempty finite gap set,
$F(U)=\max G(U)$ for its Frobenius number, and $g(U)=|G(U)|$ for its genus.
A numerical set is a \emph{numerical semigroup} if it is closed under addition.
For a proper numerical semigroup $S$, put
\[
 \PF(S)=\{f\in G(S): f+s\in S\text{ for every }s\in S\setminus\{0\}\},
 \qquad t(S)=|\PF(S)|.
\]
It is \emph{almost symmetric} if $2g(S)=F(S)+t(S)$.
A proper numerical set $U$ is \emph{symmetric} if exactly one of $a$ and
$F(U)-a$ belongs to $U$ for every integer $0\leq a\leq F(U)$.
It is \emph{reduced} when $1\notin U$, in the terminology of
S\"uer and Ye\c{s}il \cite[Definition~3.4 and Proposition~3.5]{SY}.
Their reduction deletes initial full-height columns from a Young diagram
\cite[Definition~3.2]{SY}; Proposition~3.7 describes the corresponding
shift of its numerical set. They prove that the first special subdiagram
of a symmetric numerical set is again symmetric \cite[Proposition~4.1]{SY}.
That operation differs from the row deletion considered here.

The numerical-set/Young-diagram correspondence numbers boundary steps starting
at $0$: a member gives a right step and a gap gives an up step
\cite[Section~2, pp.~4--5]{KN}.
Thus the row corresponding to a gap $r$ has length
\[
 \ell_S(r)=|S\cap\{0,\ldots,r-1\}|.
\]
Following \cite[Theorem~4.2]{DSY}, delete the rows whose \emph{original}
labels belong to $D(S)=\PF(S)\setminus\{F(S)\}$.
Every retained up step shifts left in the label sequence by the number of
deleted steps before it. Consequently the output has gap set
\begin{equation}\label{eq:gaps}
 G(\Phi(S))=
 \bigl\{r-|\{d\in D(S):d<r\}|:\ r\in G(S)\setminus D(S)\bigr\}.
\end{equation}
This formula makes sense for every proper numerical semigroup; we use $\Phi$
also for this extension, without asserting that its output is symmetric.
The Frobenius row is retained, and retained labels stay positive and strictly
increasing: between two retained labels $r<r'$ there are at most $r'-r-1$
deleted labels. Hence the formula defines a proper numerical set.
Since exactly $t(S)-1$ rows below the Frobenius row are deleted,
\begin{equation}\label{eq:shifts}
 g(\Phi(S))=g(S)-t(S)+1,\qquad
 F(\Phi(S))=F(S)-t(S)+1.
\end{equation}
On almost symmetric semigroups it is the original map, whose symmetric-output
property is proved in \cite[Theorem~4.2]{DSY}. Conversely, if the extended
output is symmetric, $2(g(S)-t(S)+1)=F(S)-t(S)+2$ forces
$2g(S)=F(S)+t(S)$. Thus a symmetric target cannot acquire a preimage
from a source that is not almost symmetric.
All intervals below contain integers unless they describe cell coordinates.

\section{The obstruction}
\begin{theorem}\label{thm:one}
For every proper numerical semigroup $S$, one has $1\notin\Phi(S)$.
In particular, $\Phi:\AS\longrightarrow\{\text{symmetric numerical sets}\}$
is not surjective, where $\AS$ denotes the proper almost symmetric
numerical semigroups.
\end{theorem}
\begin{proof}
First, $1\notin S$, since additive closure would otherwise imply $S=\N$.
If $1\notin\PF(S)$, the gap labelled $1$ is retained. No positive gap precedes
it, so its new label in \eqref{eq:gaps} is still $1$.

Suppose instead that $1\in\PF(S)$, and let $m=\min(S\setminus\{0\})$.
The defining property of $\PF(S)$ gives $s+1\in S$ for every positive
$s\in S$. Starting from $m\in S$, induction gives every integer at least $m$.
By minimality of $m$,
\[
 S=\{0\}\cup\{n\in\N:n\geq m\},\qquad
 G(S)=\{1,\ldots,m-1\},\qquad m\geq2.
\]
Every gap is then pseudo-Frobenius, because adding any positive member of $S$
gives an integer at least $m$. The deletion removes the labels
$1,\ldots,m-2$, leaving only $F(S)=m-1$, whose new label is
$(m-1)-(m-2)=1$. This includes the boundary case $m=2$, when no row is deleted.
Thus $1$ is a gap of $\Phi(S)$ in both cases.

Now let $T=\{0,1\}\cup\{n\in\N:n\geq4\}$.
It has gaps $\{2,3\}$ and Frobenius number $3$. In the complementary pairs
$\{0,3\}$ and $\{1,2\}$ exactly one element belongs to $T$, so $T$ is symmetric.
Since $1\in T$, the first assertion excludes $T$ from the image of $\Phi$.
\end{proof}

The total count in the next statement is due to Marzuola and Miller
\cite[p.~12]{MM}; the exclusion bound follows by applying Theorem~\ref{thm:one}.
\begin{corollary}\label{cor:half}
For each integer $h\geq2$, there are $2^{h-1}$ symmetric numerical sets
with Frobenius number $2h-1$, and at least $2^{h-2}$ of them are outside
the image of $\Phi$, regardless of the source Frobenius number.
\end{corollary}
\begin{proof}
The pair $\{0,2h-1\}$ has prescribed member $0$ and prescribed gap $2h-1$.
In each of the remaining $h-1$ pairs $\{a,2h-1-a\}$, $1\leq a\leq h-1$,
choose independently which element belongs to the numerical set, and include
all integers larger than $2h-1$. These choices give exactly the
$2^{h-1}$ symmetric numerical sets in question.
Exactly $2^{h-2}$ choices include $1$, and Theorem~\ref{thm:one} excludes all
of them. No proper symmetric numerical set has even Frobenius number, since
its midpoint could not satisfy the defining condition. At Frobenius number
$1$, the sole target is $\N\setminus\{1\}$, which is fixed by $\Phi$.
\end{proof}

\begin{corollary}\label{cor:fixed}
Every proper symmetric numerical semigroup $S$ is almost symmetric and
satisfies $\Phi(S)=S$. Thus every symmetric numerical semigroup occurs in
the image. Conversely, among proper numerical semigroups, these are exactly
the fixed points of the extended map.
\end{corollary}
\begin{proof}
For a gap $f<F(S)$, symmetry gives $F(S)-f\in S\setminus\{0\}$.
Since $f+(F(S)-f)=F(S)$ is a gap, $f$ is not pseudo-Frobenius.
The largest gap $F(S)$ is always pseudo-Frobenius, so
$\PF(S)=\{F(S)\}$. Symmetry pairs each gap with a member in
$\{0,\ldots,F(S)\}$, giving $2g(S)=F(S)+1=F(S)+t(S)$.
Thus $S$ is almost symmetric, and $D(S)$ is empty, so no row is deleted.

Conversely, if $\Phi(S)=S$, the genus identity in \eqref{eq:shifts} gives
$t(S)=1$, so $\PF(S)=\{F(S)\}$. For a gap $a$, choose the largest gap
$b$ in the nonempty finite set $(a+S)\cap G(S)$. Adding any positive
$s\in S$ cannot give another gap, by maximality of $b$, so $b\in\PF(S)$.
Thus $b=F(S)$ and $F(S)-a\in S$.
If instead $a\in S$ with $0\leq a\leq F(S)$, additive closure precludes
$F(S)-a\in S$. This proves symmetry.
\end{proof}

\begin{corollary}\label{cor:fiber}
Let $U_1=\N\setminus\{1\}$. Among all proper numerical semigroups,
\[
 \Phi(S)=U_1\quad\Longleftrightarrow\quad
 S=O_m:=\{0\}\cup\{n\in\N:n\geq m\}\quad\text{for some }m\geq2.
\]
Every $O_m$ is almost symmetric. Hence $U_1$ has an infinite fiber, with
exactly one source of each positive Frobenius number.
\end{corollary}
\begin{proof}
If $\Phi(S)=U_1$, \eqref{eq:shifts} gives $g(S)=t(S)$.
All gaps are then pseudo-Frobenius; in particular $1\in\PF(S)$.
The proof of Theorem~\ref{thm:one} shows that $S=O_m$.
Conversely, that proof gives $\Phi(O_m)=U_1$.
For $O_m$, $F=m-1$ and $g=t=m-1$, so $2g=F+t$.
The only ordinary source of Frobenius $F$ is $O_{F+1}$.
\end{proof}

\begin{corollary}\label{cor:source-bound}
For $F\geq1$, let $\AS_F$ be the almost symmetric numerical semigroups
of Frobenius number $F$, $\mathcal S_F=\Phi(\AS_F)$, and
$\mathcal{NS}_F$ the semigroups in $\mathcal S_F$, as in
\cite[Section~5, pp.~8--9]{DSY}. Then
\[
 1\leq|\mathcal{NS}_F|\leq|\mathcal S_F|
 \leq 2^{\lfloor(F+1)/2\rfloor-1}.
\]
Every output has Frobenius number $2(F-g(S)+1)-1\leq F$.
\end{corollary}
\begin{proof}
For $S\in\AS_F$, almost symmetry and \eqref{eq:shifts} give
$g(\Phi(S))=F-g(S)+1$ and $F(\Phi(S))=2g(\Phi(S))-1$.
Since $t(S)\geq1$, the latter is at most $F$.
Put $H=\lfloor(F+1)/2\rfloor$. At output Frobenius number $1$ there is
one symmetric target. At output Frobenius number $2h-1$, $h\geq2$,
Theorem~\ref{thm:one} leaves at most $2^{h-2}$ possible targets,
by the pair count in Corollary~\ref{cor:half}.
Summing over $1\leq h\leq H$ gives
$1+\sum_{h=2}^{H}2^{h-2}=2^{H-1}$, also when $H=1$.
Finally $O_{F+1}\in\AS_F$ has semigroup image $U_1$, proving the lower bound.
\end{proof}

\section{The prefix obstruction}
Restricting the target to reduced symmetric numerical sets does not restore
surjectivity. For numerical semigroups themselves, the square-diagram path
is a Dyck path \cite[Theorem~4.3 in the arXiv version]{BM}; the equivalent
gap and ballot conditions are also stated in
\cite[Theorem~7.2 and Remark~7.3]{RR}.
Here $\Phi(S)$ need not be a semigroup, so the output prefix inequality
requires the retained-label argument below.
\begin{theorem}\label{thm:bound}
Let $S$ be a proper almost symmetric numerical semigroup and let
$T=\Phi(S)$. Write $F(T)=2h-1$ and
$T\cap[0,F(T)]=\{0=b_0<b_1<\cdots<b_{h-1}\}$.
Then $b_i\geq2i$ for every $0\leq i<h$. Equivalently,
\[
 |T\cap[0,r]|\leq\lfloor r/2\rfloor+1
 \qquad(0\leq r\leq F(T)).
\]
\end{theorem}
\begin{proof}
Deleting $t-1$ original labels below $F$ gives $F(T)=F-t+1$.
Almost symmetry therefore gives
$h=(F-t+2)/2=F+1-g=|S\cap[0,F]|$.
Put $A=S\cap[0,F]=\{0=a_0<a_1<\cdots<a_{h-1}\}$.
We first identify all retained labels. Closure under addition implies
$F-A\subseteq G(S)$. For $a\in A\setminus\{0\}$,
$(F-a)+a=F\notin S$, so $F-a\notin\PF(S)$.
As $F\in\PF(S)$, it follows that
$D\subseteq G(S)\setminus(F-A)$. Almost symmetry gives
\[
 |G(S)\setminus(F-A)|=g-(F+1-g)=t-1=|D|.
\]
Thus $D=G(S)\setminus(F-A)$, also the identity
in~\cite[Proposition~3.2(3)]{DSY}; compare also the almost-symmetry
identity in~\cite[proof of Proposition~2.9]{CKLOS}. The retained labels are exactly
the disjoint union $R=A\sqcup(F-A)$.
Deletion relabels a retained $u$ by its rank
\[
 c(u)=|\{v\in R:v<u\}|=u-|\{d\in D:d<u\}|.
\]
There are $2h$ retained labels and the last is the gap $F$.
Hence $F(T)=2h-1$ and $b_i=c(a_i)$.
Let $q_i=|\{a\in A:a+a_i>F\}|$. Exactly $i$ elements of $A$
precede $a_i$, and exactly $q_i$ members of $F-A$ precede it. Therefore
$b_i=i+q_i$. Translation by $a_i$ gives an injection
\[
 A\cap[0,F-a_i]\longrightarrow A\cap[a_i,F],\qquad
 a\longmapsto a+a_i,
\]
since $S$ is closed under addition. The two sets have sizes $h-q_i$
and $h-i$, respectively, so $q_i\geq i$ and $b_i\geq2i$.
If $\ell=|T\cap[0,r]|$, then $2(\ell-1)\leq b_{\ell-1}\leq r$,
giving the equivalent prefix bound. Conversely, applying that bound
at $r=b_i$ yields $b_i\geq2i$.
\end{proof}

\begin{corollary}\label{cor:example}
The reduced symmetric numerical set
$T_* = \{0,2,3,6\}\cup\{8,9,\ldots\}$ is not in the image of $\Phi$.
Its Frobenius number $7$ is minimal for such a counterexample.
\end{corollary}
\begin{proof}
Its gaps are $\{1,4,5,7\}$, and exactly one member of each pair
$\{0,7\},\{1,6\},\{2,5\},\{3,4\}$ lies in $T_*$.
It is reduced and symmetric, but $b_2=3<4$,
contrary to Theorem~\ref{thm:bound}.
Every proper symmetric numerical set has odd Frobenius number, since an even
Frobenius number would pair its midpoint with itself.
For Frobenius numbers $1,3,5$, the reduced symmetric sets are respectively
\[
 \{0\}\cup[2,\infty),\qquad
 \{0,2\}\cup[4,\infty),\qquad
 \{0,2,4\}\cup[6,\infty),\quad
 \{0,3,4\}\cup[6,\infty),
\]
where intervals denote integers. Each is a symmetric numerical semigroup.
For a symmetric numerical semigroup $\PF(S)=\{F(S)\}$,
so $D$ is empty and $\Phi(S)=S$.
\end{proof}

Taking $r=h-1$ in Theorem~\ref{thm:bound} shows that
$k=|T\cap[1,h-1]|$ satisfies $2k+1\leq h$.
For $h\geq2$, the reduced symmetric targets of Frobenius number $2h-1$
are parametrized by their elements in $\{2,\ldots,h-1\}$,
so there are $2^{h-2}$ targets. Theorem~\ref{thm:bound} excludes at least
\[
 \sum_{j=\lceil h/2\rceil}^{h-2}\binom{h-2}{j}
\]
of them, with an empty sum interpreted as zero.

There are also excluded reduced symmetric targets with every least positive
element $m\geq2$:
\[
 T_m=\{0\}\cup[m,2m-1]\cup[3m,4m-2]\cup[4m,\infty).
\]
The map $x\mapsto F(T_m)-x$, with $F(T_m)=4m-1$, exchanges $0$ with $4m-1$,
$[1,m-1]$ with $[3m,4m-2]$, and $[m,2m-1]$ with $[2m,3m-1]$.
Thus $T_m$ is symmetric and reduced. However,
$|T_m\cap[0,2m-1]|=m+1$, whereas the prefix bound permits at most $m$.

\section{Rigidity for targets containing \texorpdfstring{$2$}{2}}
\begin{theorem}\label{thm:two}
Let $T$ be a reduced symmetric numerical set with $2\in T$ and
$F(T)=2h-1$. Then $T$ is in the image of $\Phi$ if and only if
\[
 T=\{0,2,4,\ldots,2h-2\}\cup[2h,\infty).
\]
\end{theorem}
\begin{proof}
The case $h=1$ is the single set $\{0\}\cup[2,\infty)$.
Suppose $h\geq2$ and $T=\Phi(S)$, and use the notation of
Theorem~\ref{thm:bound}. Put $\mu=a_1$. Since $b_1=2$, the identity
$b_1=1+q_1$ gives $q_1=1$. Translation by $\mu$ is therefore a bijection
\[
 A\setminus\{\max A\}\longrightarrow A\setminus\{0\}.
\]
It preserves order, so $a_{i+1}=a_i+\mu$ and
$A=\{0,\mu,\ldots,(h-1)\mu\}$.
Now $(h-1)\mu<F<h\mu$: the first inequality holds because $F$ is a gap,
and the second because $\max A+\mu\notin A$ and belongs to $S$.
Write $F=(h-1)\mu+r$ with $1\leq r<\mu$.
The retained elements and gaps occur as the successive pairs
\[
 (j\mu,\ j\mu+r),\qquad j=0,\ldots,h-1,
\]
in that order, since $0<r<\mu$.
Thus every retained source element has an even rank, giving the displayed
form of $T$. Conversely, this $T$ is the symmetric numerical semigroup
$\langle2,2h+1\rangle$, so $\Phi(T)=T$.
\end{proof}

The prefix inequalities alone are not sufficient. The set
$U=\{0,2,5,6,8\}\cup[10,\infty)$ is reduced and symmetric, and its
elements below the conductor, $(0,2,5,6,8)$, satisfy $b_i\geq2i$.
This symmetric numerical set already appears in~\cite[Example~5.3]{Y}.
Here $2\in U$ and $4\notin U$, so Theorem~\ref{thm:two} establishes
the additional conclusion that $U$ is outside the row-deletion image.


\section{The full-cell restriction}\label{sec:cells}
For a symmetric numerical set $T$, write $F(T)=2h-1$ and
$T\cap[0,2h-1]=\{0=b_0<\cdots<b_{h-1}\}$.
Use the Young diagram convention of Delgado, S\"uer and
Ye\c{s}il~\cite[Section~2, p.~4]{DSY}: read $0,\ldots,F(T)$,
taking an east step for an element and a north step for a gap.
The diagram is the region above this path and below height $h$.
We forbid a full cell below $y=x$ when its entire interior is in $y<x$;
cells bisected by the diagonal are allowed. A cell touching the diagonal
at its northwest corner, with all its interior below, is forbidden.

\begin{lemma}\label{lem:cells}
The full-cell condition is equivalent to $b_i\geq2i$ for every $i$.
\end{lemma}
\begin{proof}
The east step labelled $b_i$ starts at $(i,b_i-i)$.
Column $i$ therefore contains precisely the cells
$[i,i+1]\times[j,j+1]$ with $b_i-i\leq j\leq h-1$.
Such a cell has its entire interior below $y=x$ exactly when $j+1\leq i$.
There is no such cell in any column if and only if $b_i-i\geq i$.
\end{proof}
This is not the condition that the boundary path itself stays above $y=x$:
its initial east step is allowed. Numerical-semigroup boundary paths and
their gap inequalities are classical~\cite[Theorem~4.3]{BM}; here the
condition is imposed on the symmetric output set, which need not be closed
under addition.

\Needspace{11\baselineskip}
\section{Excluded targets on the square board}
\begin{theorem}\label{thm:diagonal}
The reduced symmetric numerical set
\[
 U=\{0,2,5,6,8\}\cup[10,\infty)
\]
satisfies the full-cell condition but has no almost symmetric numerical
semigroup preimage. With the convention in Section~\ref{sec:cells}, $9$ is the smallest
Frobenius number of such an excluded target. For every $h\geq5$,
\[
 U_h=\{0,2\}\cup\bigl([h,2h-2]\setminus\{2h-3\}\bigr)\cup[2h,\infty)
\]
gives a family of such targets with Frobenius number $2h-1$, with $U_5=U$.
\end{theorem}
\begin{proof}
For $U$, the complementary pairs are $\{0,9\}$, $\{1,8\}$,
$\{2,7\}$, $\{3,6\}$ and $\{4,5\}$; each contains exactly one element
of $U$. Its elements $(0,2,5,6,8)$ satisfy $b_i\geq2i$.
Yet $2\in U$ and $4\notin U$, contradicting Theorem~\ref{thm:two}
if $U$ were an image. The numerical set itself is depicted by
Bras-Amor\'os and de Mier~\cite[Remark~2.5]{BM} and appears explicitly
in Ye\c{s}il~\cite[Example~5.3]{Y}; the point here is its exclusion
despite the diagonal restriction.

For $U_h$, the low labels $0,2$ are paired with the omitted high labels
$2h-1,2h-3$. Each other low label in $[0,h-1]$ is omitted and has its
complement in the displayed high interval. Thus $U_h$ is symmetric,
reduced, and has Frobenius number $2h-1$. Its ordered elements below the
conductor are
\[
 b_0=0,\quad b_1=2,\quad
 b_i=h+i-2\ (2\leq i\leq h-2),\quad b_{h-1}=2h-2.
\]
These satisfy $b_i\geq2i$, while $2\in U_h$ and $4\notin U_h$.
The same rigidity argument excludes it.

For minimality, symmetry and $b_i\geq2i$ give exactly the following lists
below the conductor when $h\leq4$. They follow by choosing one member
of each pair $\{j,2h-1-j\}$, retaining $0$, and imposing the inequalities.
\begin{center}
\begin{tabular}{cl}
\toprule
$F(T)$ & Possible $T\cap[0,F(T)]$\\
\midrule
1 & $\{0\}$\\
3 & $\{0,2\}$\\
5 & $\{0,2,4\}$, $\{0,3,4\}$\\
7 & $\{0,2,4,6\}$, $\{0,3,5,6\}$, $\{0,4,5,6\}$\\
\bottomrule
\end{tabular}
\end{center}
Each completed set is closed under addition, hence is a symmetric
numerical semigroup fixed by $\Phi$. This proves the minimality assertion.
\end{proof}

\begin{figure}[ht]
\centering
\setlength{\unitlength}{6mm}
\begin{picture}(7,5.6)
\put(0,0){\framebox(1,1){}}
\multiput(0,1)(1,0){2}{\framebox(1,1){}}
\multiput(0,2)(1,0){2}{\framebox(1,1){}}
\multiput(0,3)(1,0){4}{\framebox(1,1){}}
\multiput(0,4)(1,0){5}{\framebox(1,1){}}
\linethickness{0.8pt}\put(0,0){\line(1,1){5}}
\put(5.2,4.9){\makebox(0,0)[l]{$y=x$}}
\put(0,-0.45){\makebox(0,0){$0$}}
\put(5,-0.45){\makebox(0,0){$5$}}
\end{picture}
\caption{The diagram of $U$, with row lengths $1,2,2,4,5$ from bottom
to top. The diagonal may bisect a cell; no full cell lies below it.}
\end{figure}

The exclusion and family statements have written proofs for every source
and every $h\geq5$. The accompanying \texttt{verify\_diagonal.py} checks all symmetric
targets through Frobenius number $23$, the family for $5\leq h\leq100$,
and all numerical-semigroup sources through Frobenius number $18$.


\section{Multiplicity three and a general membership bound}\label{sec:three}
The multiplicity of a numerical set is its least positive element.
We thank Mehmet Ye\c{s}il for suggesting the question at multiplicity three.
For a symmetric target $T$ with $F(T)=2h-1$, write
$B=T\cap[0,2h-1]$; complete every displayed $B$ by adjoining $[2h,\infty)$.
For $h\geq3$, define
\begin{align}
 E_h&=\{0\}\cup\{3,5,\ldots,2h-3\}\cup\{2h-2\},\label{eq:E}\\
 P_h&=\langle3,h+1\rangle\cap[0,2h-1]
                 &&\text{if }h\equiv0,1\pmod3.\label{eq:P01}
\end{align}
When $h=3k+2$ with $k\geq1$, put
\begin{equation}\label{eq:P2}
 P_h=\{3j:0\leq j\leq k\}
 \cup\{3k+2+3j:0\leq j\leq k\}
 \cup\{3k+4+3j:0\leq j\leq k-1\}.
\end{equation}
Here $\langle u,v\rangle=\{\alpha u+\beta v:\alpha,\beta\in\N\}$.

\begin{theorem}\label{thm:classification}
Let $T$ be a symmetric numerical set of multiplicity three, with
$F(T)=2h-1$. Then $h\geq3$, and $T=\Phi(S)$ for a proper almost symmetric
numerical semigroup $S$ if and only if $B=E_h$ or $B=P_h$.
The families coincide for $h=3,4,5$ and are distinct for every $h\geq6$.
\end{theorem}

The small cases make the difference from multiplicity two visible:
\begin{center}
\begin{tabular}{cll}
\toprule
$h$ & $E_h$ & $P_h$\\
\midrule
3 & $0,3,4$ & $0,3,4$\\
4 & $0,3,5,6$ & $0,3,5,6$\\
5 & $0,3,5,7,8$ & $0,3,5,7,8$\\
6 & $0,3,5,7,9,10$ & $0,3,6,7,9,10$\\
8 & $0,3,5,7,9,11,13,14$ & $0,3,6,8,10,11,13,14$\\
\bottomrule
\end{tabular}
\end{center}

The proof reduces source multiplicity to three or four, then uses elementary
Ap\'ery thresholds. For related work on source multiplicities three and four,
see Rosales~\cite{Rosales}; the issue here is their common row-deletion
image. Almost symmetry and its gap reflection are treated in
Nari~\cite[Section~2]{Nari}. We include the needed arguments so that the
classification uses no finite-search premise.


\subsection{Residue chains and the source cutoff}
\begin{lemma}\label{lem:rank}
Let $S$ be a proper almost symmetric numerical semigroup, using the
notation $A,F,h,b_i$ of Theorem~\ref{thm:bound}. If $h\geq2$, write
$\mu=a_1$ for the source multiplicity and $r=b_1$ for the output multiplicity.
Then $A$ consists of exactly $r-1$ nonempty residue chains modulo $\mu$,
each with step $\mu$ and no internal missing element.
\end{lemma}
\begin{proof}
In any residue class occurring in $A$, addition by $\mu$ fills the chain
from its first element to its last element below $F$. Its last element
is in $(F-\mu,F]$, and each element in this interval is a chain end.
Thus the number of chains is $|A\cap(F-\mu,F]|=r-1$ by~the rank formula in Theorem~\ref{thm:bound}
at $i=1$.
\end{proof}

\begin{theorem}\label{thm:cutoff}
Let $S$ be a proper almost symmetric numerical semigroup. Suppose $\Phi(S)$ has multiplicity $r$ and Frobenius number $2h-1$,
with $h>r$. Then
\begin{equation}\label{eq:bound}
 \mu(S)\leq2r-2,\qquad
 2h-1\leq F(S)\leq2h+2r-5.
\end{equation}
\end{theorem}
\begin{proof}
If $F<2\mu$, the chain of residue zero has at most two elements, while
each other chain has at most one: its first element is greater than $\mu$.
Lemma~\ref{lem:rank} would give $h\leq2+(r-2)=r$. Thus $F\geq2\mu$.
All $f=1,\ldots,\mu-1$ are gaps, and $f+\mu\leq F$.
At most $r-2$ of these translates can belong to $A$, one in each nonzero
active residue class. For at least $\mu-r+1$ values of $f$, therefore,
$f+\mu\notin S$, so $f\notin\PF(S)$. The gap-reflection identity in
Theorem~\ref{thm:bound} puts their distinct complements $F-f$ in $A$.
All these complements lie in $(F-\mu,F]$, which has exactly $r-1$
elements of $A$. Hence $\mu-r+1\leq r-1$.

No pseudo-Frobenius gap is divisible by $\mu$. Two such gaps cannot
have the same nonzero residue: their positive difference is a multiple
of $\mu$ in $S$, contradicting the pseudo-Frobenius property of the
smaller gap. Thus $1\leq t\leq\mu-1$.
Since $h=F+1-g$ and $2g=F+t$, we have $F=2h+t-2$.
Combining these facts gives~\eqref{eq:bound}.
\end{proof}

Theorem~\ref{thm:one} excludes multiplicity one. For a symmetric numerical set with $h\geq2$, its multiplicity is at most
$h$: one of $h-1,h$ must be an element. If $r=h$, symmetry forces
$B=\{0\}\cup[h,2h-2]$; this completed set is a numerical semigroup
and is fixed by $\Phi$. The remaining one-gap set is
$\N\setminus\{1\}$, also fixed by $\Phi$.
For every remaining symmetric target, with $2\leq r<h$, Theorem~\ref{thm:cutoff} gives
an exact finite membership procedure: enumerate the numerical semigroups
with $2h-1\leq F\leq2h+2r-5$ and multiplicity at most $2r-2$, retain
the almost symmetric ones, and compare their compressed labels with $B$.
Necessity follows from the theorem; sufficiency follows by exhibiting a
source. This procedure decides membership, not the number of all sources.


\subsection{Proof of the classification}\label{sec:three-proof}
When $r=3$ and $h=3$, symmetry forces $B=\{0,3,4\}$, realized by
$\langle3,4\rangle$. Assume henceforth that $h\geq4$.
Theorem~\ref{thm:cutoff} gives $\mu\leq4$. Multiplicity two is impossible:
its even chain already contains zero and the sole odd residue class
can first occur only beyond $F$, since otherwise adding two would put the
odd gap $F$ in $S$. Thus it has one active chain, not two.

\subsubsection{Source multiplicity four}
There are two active residues, $0$ and $\rho\in\{1,2,3\}$.
Since $h>3$, the proof of Theorem~\ref{thm:cutoff} gives $F\geq8$;
as $8\in S$, actually $F\geq9$.
For $f\in\{1,2,3\}\setminus\{\rho\}$, the integer $f+4$ is a gap.
Hence the two complements $F-f$ lie in the two active residue classes.
If $4+\rho$ were also a gap, all three complements would belong to
$A\cap(F-4,F]$, contradicting its cardinality two. Thus $4+\rho\in A$.
The complement residues now give:
\begin{itemize}
\item If $\rho=1$, then $F\equiv3\pmod4$. The complement $F-1$ has
inactive residue two, so $1\in\PF(S)$. But $5\in S$ and $1+5=6\notin S$,
a contradiction.
\item If $\rho=2$, then $F$ is odd and $6\in S$. The two chains are
precisely $A=\{0\}\cup\{4,6,\ldots,F-1\}$. Their cardinality gives
$F=2h+1$.
\item If $\rho=3$, then $F\equiv1\pmod4$, and $3\in\PF(S)$ because
$F-3$ has inactive residue two. If $10\leq F$, its inactive residue
makes it a gap. But $3\in\PF(S)$ and $7\in S$ imply $10\in S$.
Hence $10>F$. Thus $F=9$, $A=\{0,4,7,8\}$ and $h=4$.
\end{itemize}
In the second case the source is
\begin{equation}\label{eq:even-source}
 S_h=\{0\}\cup\{4,6,\ldots,2h\}\cup[2h+2,\infty).
\end{equation}
Its pseudo-Frobenius gaps are exactly $2,2h-1,2h+1$: each earlier odd
gap fails on addition by four, while the displayed three satisfy the
definition. Its genus is $h+2$, so it is almost symmetric of type three.
Deleting $2$ and $2h-1$ sends its elements below $F$ to $E_h$.
This also proves existence of $E_h$ for every $h\geq3$.
In the last case the deleted gaps are $3,6$, giving $E_4=P_4$.

\subsubsection{Source multiplicity three}
Let the first elements in residues one and two be
$w_1=3\alpha+1$ and $w_2=3\beta+2$, where $\alpha,\beta\geq1$.
These are the nonzero elements of $\Ap(S,3)$.
Addition gives $\beta\leq2\alpha$ and $\alpha\leq2\beta+1$.
Counting gaps in these two residues yields
\[
 g=\alpha+\beta,\qquad
 F=\max(3\alpha-2,3\beta-1).
\]
The residue argument in Theorem~\ref{thm:cutoff} gives $t\in\{1,2\}$.
Substitution in $2g=F+t$ leaves four cases:
\begin{center}
\begin{tabular}{cccc}
\toprule
$F$ & $t$ & Threshold relation & $h$\\
\midrule
$3\alpha-2$ & 1 & $\alpha=2\beta+1$ & $3\beta+1$\\
$3\beta-1$ & 1 & $\beta=2\alpha$ & $3\alpha$\\
$3\alpha-2$ & 2 & $\alpha=2\beta$ & $3\beta-1$\\
$3\beta-1$ & 2 & $\beta=2\alpha-1$ & $3\alpha-2$\\
\bottomrule
\end{tabular}
\end{center}
In the first two rows the larger threshold is twice the smaller one.
Thus $S=\langle3,h+1\rangle$, a symmetric semigroup, and nothing is
deleted. Its output is~\eqref{eq:P01}.

In either type-two row there is exactly one deleted gap. The retained
labels $A\sqcup(F-A)$ are reflection-invariant, so the singleton deleted
set is also reflection-invariant. It is therefore $\{F/2\}=\{h\}$.
Every source element greater than $h$ decreases by one; all others stay.
In the third row write $\beta=k+1$, so $h=3k+2$ and $k\geq1$.
Here
\[
 A=\{3j:0\leq j\leq2k+1\}
   \cup\{3k+5+3j:0\leq j\leq k-1\}.
\]
Decreasing exactly the elements above $3k+2$ gives~\eqref{eq:P2}.
These sources exist for every $k\geq1$: the semigroup
$\langle3,3k+5,6k+7\rangle$ has Ap\'ery set
$\{0,6k+7,3k+5\}$ and pseudo-Frobenius gaps $3k+2,6k+4$.
Indeed adding three to either gap gives its threshold, and adding either
nonzero threshold gives an element; all other gaps fail on addition by
three. Thus its type is two and $2g=F+2$.

In the fourth row put $\alpha=a\geq2$, giving $h=3a-2$.
The elements of $A$ before deletion are the multiples of three through $6a-6$
and $3a+1,3a+4,\ldots,6a-5$. After deleting the gap $3a-2$, they become
\[
 \{3j:0\leq j\leq2a-2\}
 \cup\{3a-1+3j:0\leq j\leq a-2\}.
\]
This is exactly $\langle3,3a-1\rangle\cap[0,6a-5]$, again $P_h$.
For $h\geq4$, all possible sources have now been exhausted, and the constructions give
both families. The lists for $h=3,4,5$ agree. For $h\geq6$, $6\in P_h$
but $6\notin E_h$, proving distinctness and completing
Theorem~\ref{thm:classification}.


\subsection{An excluded target of multiplicity three}\label{sec:three-example}
The symmetric set
\[
 V=\{0,3,4,7,8\}\cup[10,\infty)
\]
has multiplicity three but is neither $E_5$ nor $P_5$, and hence has no
almost symmetric source. Writing $V\cap[0,9]=\{b_0,\ldots,b_4\}$ in
increasing order, we have $b_i\geq2i$ for $0\leq i\leq4$.
Thus those rank inequalities remain insufficient at multiplicity three.
Frobenius number nine is minimal for this exclusion: at Frobenius numbers
five and seven, the only multiplicity-three symmetric sets have left
elements $\{0,3,4\}$ and $\{0,3,5,6\}$, respectively.


\section{Finite verification}
The accompanying program \texttt{verify\_prefix.py} enumerates all $840\,170$ proper
numerical semigroups with $F(S)\leq35$. Among them $17\,842$ are
almost symmetric. For each it computes $\PF(S)$ from its defining
condition, deletes the original rows, and checks symmetry and
Theorems~\ref{thm:bound} and~\ref{thm:two}. None has image $T_*$.
The image counts below are unions over this bounded \emph{source}
range, grouped by \emph{output} Frobenius number; they are not counts
for sources of one fixed Frobenius number.
\begin{center}
\setlength{\tabcolsep}{3pt}
\begin{tabular}{lrrrrrrrrrrr}
\toprule
$F(T)$&1&3&5&7&9&11&13&15&17&19&21\\
Reduced targets&1&1&2&4&8&16&32&64&128&256&512\\
Attained, $F(S)\leq35$&1&1&2&3&4&7&9&12&18&25&32\\
\bottomrule
\end{tabular}
\end{center}
To see completeness of the enumeration, adjoining the Frobenius gap
of any proper numerical semigroup gives its unique parent semigroup.
Conversely, a child removes a minimal positive generator larger than
the parent's Frobenius number. For parent multiplicity $m$, all such
generators are at most $F+m$: a larger integer $x$ decomposes as
$m+(x-m)$, with $x-m>F$. Starting at $\N\setminus\{1\}$ and keeping
removals at most $35$ therefore visits every proper source in the range,
exactly once. For a removed label $x$, minimality is tested by all
decompositions $x=a+(x-a)$ with $1\leq a\leq x/2$.

For the row-deletion formula extended to any proper numerical semigroup,
the output has genus $g-t+1$ and Frobenius number $F-t+1$.
If that output is symmetric, their relation forces $2g=F+t$.
Thus a non-almost-symmetric source cannot produce $T_*$ either.
Theorems~\ref{thm:bound} and~\ref{thm:two}, and Corollary~\ref{cor:example}, have written
proofs for all parameters; the enumeration is finite supporting evidence.

The program \texttt{verify\_multiplicity3.py} checks the two classified
families, their explicit sources for $3\leq h\leq100$, the general cutoff,
and the excluded example. Through source Frobenius $35$ there are $95$
source instances with multiplicity-three outputs, counting repeated targets,
and $11\,448$ instances of the general bound. The program independently
enumerates all $1\,654$ proper sources through Frobenius $18$ by testing
subsets for closure. For multiplicity three, the cutoff makes the image
lists complete through output Frobenius $33$, with the $h=3$ boundary
settled in Theorem~\ref{thm:classification}.

\section{Relation to the neighboring questions}
The fixed-source bound concerns the paper's $\mathcal S_F$, not all targets
of a fixed output Frobenius number. The exact ordinary fiber strengthens
the source's noninjectivity examples. The bound does not decide the asymptotic ratios plotted in Section~5. Theorems~\ref{thm:classification} and~\ref{thm:cutoff} give a closed classification at multiplicity three and a finite membership test in general, but not closed families at higher multiplicity.

\smallskip
\begingroup\small
\noindent\begin{tabularx}{\linewidth}{@{}p{0.40\linewidth}X@{}}
\toprule
Source item in \cite{DSY} & Coverage \\
\midrule
Section~3, Proposition~3.2; Section~4, Theorem~4.2 & Prior results cited:
characterization and symmetric-output construction of Delgado, S\"uer and Ye\c{s}il. \\
Section~5, p.~8, noninjectivity & Prior proof cited; Corollary~\ref{cor:fiber}
identifies an entire infinite fiber. \\
Section~5, p.~9, final paragraph: surjectivity & Answered: no
(Theorem~\ref{thm:one}); half of the targets at each odd Frobenius number
at least $3$ are excluded (Corollary~\ref{cor:half}). \\
Section~5, pp.~8--9: how many outputs are semigroups? & Partial answer:
all symmetric semigroups are attained (Corollary~\ref{cor:fixed});
fixed-source bounds in Corollary~\ref{cor:source-bound}; no exact general count. \\
Section~5, p.~9: image and plotted ratios & One full fiber is determined;
the prefix and full-cell conditions are necessary, not sufficient (Theorems~\ref{thm:bound} and~\ref{thm:diagonal}); closed families at higher multiplicity and ratio asymptotics are not covered. \\
Ye\c{s}il's earlier consecutive-PF construction \cite[Lemma~5.7]{Y} &
Prior symmetric-output proof cited; the omission of $1$ also holds for this
restricted construction, by Theorem~\ref{thm:one}. \\
Fixed points and source-count bound & Not posed in the source;
proved in Corollaries~\ref{cor:fixed} and~\ref{cor:source-bound}. \\
Total symmetric-target count & Prior proof: Marzuola--Miller
\cite[p.~12]{MM}; the image exclusion is proved in Corollary~\ref{cor:half}. \\
Reduced targets and the diagonal restriction &
Corollary~\ref{cor:example} gives the minimal reduced nonimage;
Theorem~\ref{thm:two} characterizes images containing $2$;
Lemma~\ref{lem:cells} identifies the full-cell convention and
Theorem~\ref{thm:diagonal} gives the minimal full-cell nonimage and a family. \\
Targets of multiplicity three; general membership &
Theorem~\ref{thm:classification} and Section~\ref{sec:three-example} give the exact images and a minimal prefix-compatible nonimage;
Theorem~\ref{thm:cutoff} supplies the cutoff for the general membership test. \\
\bottomrule
\end{tabularx}
\endgroup

\paragraph{Formal statements and finite checks.}
In the pinned \href{https://github.com/the-omega-institute/trureturing/blob/f960cb09514c99f42fd07398f5240a77378ea4db/D5/S3/Combinatorics/Partitions/DelgadoSuerYesilPhiRefutation.lean}{Lean source file},
the declaration \texttt{phi\_omits\_one} states the universal obstruction for
proper numerical semigroups, and \texttt{result} states the negation of
unrestricted surjectivity, using the target with gaps $\{2,3\}$.
The file's \texttt{deleteRows} and \texttt{phi} definitions implement
\eqref{eq:gaps}; its \texttt{AlmostSymmetric} definition is additive closure
together with $2g=F+t$. The namespace is
\texttt{D5.S3.Combinatorics.Partitions.}\allowbreak
\texttt{DelgadoSuerYesilPhiRefutation}.
The four corollaries in Section~2, the prefix and rigidity theorems,
the full-cell equivalence, the two minimality and family arguments,
and the residue-chain, membership and multiplicity-three results
have written proofs above; they are separate from these two formal declarations.
The accompanying \texttt{verify\_base.py} uses exact integer and set operations to
check the row-deletion identities, obstruction, ordinary fiber and fixed-point
equivalence for every proper numerical semigroup with $1\leq F\leq18$,
the fixed-source bound in that range, and all symmetric targets with odd
$1\leq F\leq17$. The precise counts and Table~1 comparisons are supplied
with the checker. These finite checks supplement the unbounded proofs.


The independent membership-mask implementation
\texttt{verify\_independent.py} checks the prefix and rigidity assertions
through source Frobenius number $16$ and agrees with the semigroup-tree
counts throughout that range. The full-cell equivalence, small-target lists
and excluded family are also checked by \texttt{verify\_diagonal.py} in
the ranges stated above.

\begin{thebibliography}{99}
\small
\bibitem{DSY}
M.~Delgado, M.~S\"uer and M.~Ye\c{s}il,
\emph{A Symmetrization for Almost Symmetric Numerical Semigroups},
arXiv:2611.14268v1 (2026),
\href{https://arxiv.org/abs/2611.14268v1}{arxiv.org/abs/2611.14268v1}.

\bibitem{KN}
W.~J.~Keith and R.~Nath,
\emph{Partitions with prescribed hooksets},
J.~Combinatorics Number Theory \textbf{3} (2011), no.~1, 39--50.
Author version: arXiv:1011.1945v2 (2010), Section~2, pp.~4--5,
\href{https://arxiv.org/abs/1011.1945v2}{arxiv.org/abs/1011.1945v2}.

\bibitem{Y}
M.~Ye\c{s}il,
\emph{Young Diagram Decompositions for Almost Symmetric Numerical Semigroups},
Results Math.~\textbf{80} (2025), article~217,
\href{https://doi.org/10.1007/s00025-025-02527-x}{doi:10.1007/s00025-025-02527-x}.
Author version: arXiv:2504.20688v2 (2026), Lemma~5.7,
\href{https://arxiv.org/abs/2504.20688v2}{arxiv.org/abs/2504.20688v2}.

\bibitem{SY}
M.~S\"uer and M.~Ye\c{s}il,
\emph{Special subdiagrams of Young diagrams and numerical semigroups},
Turkish J.~Math.~\textbf{48} (2024), no.~2, 346--359,
\href{https://doi.org/10.55730/1300-0098.3510}{doi:10.55730/1300-0098.3510}.

\bibitem{MM}
J.~Marzuola and A.~Miller,
\emph{Counting numerical sets with no small atoms}, arXiv:0805.3493v1 (2008),
p.~12, section ``Symmetric numerical sets with no small atoms'',
\href{https://arxiv.org/abs/0805.3493v1}{arxiv.org/abs/0805.3493v1}.

\bibitem{BM}
M.~Bras-Amor\'os and A.~de Mier,
\emph{Representation of Numerical Semigroups by Dyck Paths},
Semigroup Forum \textbf{75} (2007), 676--681.
\url{https://doi.org/10.1007/s00233-007-0717-7}.
Theorem~4.3 in the preprint version:
\url{https://arxiv.org/abs/math/0612634v1}.

\bibitem{RR}
R.~Raja and Z.~Rashid,
\emph{Creation Sequences for Numerical Semigroups and Exact Counts of
Semigroup Partitions}, arXiv:2611.20832v1 (2026), Theorem~7.2 and Remark~7.3.
\url{https://arxiv.org/abs/2611.20832v1}.

\bibitem{CKLOS}
A.~Chen, N.~Kaplan, L.~Lawson, C.~O'Neill and D.~Singhal,
\emph{Enumerating numerical sets associated to a numerical semigroup},
arXiv:2211.17090v2 (2023), Proposition~2.9.
\url{https://arxiv.org/abs/2211.17090v2}.

\bibitem{MZ}
H.~Ma and W.~Zhang,
\emph{A negative answer to Delgado, S\"uer and Ye\c{s}il's symmetrization question},
revision~2 (2026), Section~2.
\url{https://doi.org/10.5281/zenodo.24391331}.
\bibitem{Nari} H.~Nari, \emph{Symmetries on almost symmetric numerical
semigroups}, Semigroup Forum \textbf{86} (2013), 140--154.
\url{https://doi.org/10.1007/s00233-012-9397-z}.
\bibitem{Rosales} J.~C.~Rosales, \emph{Numerical Semigroups with
Multiplicity Three and Four}, Semigroup Forum \textbf{71} (2005), 323--331.
\url{https://doi.org/10.1007/s00233-005-0517-x}.

\bibitem{MZ4} H.~Ma and W.~Zhang, \emph{Extensions of Ma--Zhang's negative answer
to Delgado, S\"uer and Ye\c{s}il's symmetrization question}, revision 4 (2026),
Sections~3--6. \url{https://doi.org/10.5281/zenodo.24507291}.
\end{thebibliography}

\smallskip
\noindent\small Copyright \copyright\ 2026 Haobo Ma and Wenlin Zhang.
Licensed under Creative Commons Attribution 4.0 International (CC BY 4.0).
\end{document}
